\documentclass[../principal.tex]{subfiles}
\graphicspath{{\subfix{../imagenes/}}}

\begin{document}



\chapter{Variables}

En computación usamos variables para almacenar información. Existen muchos tipos, dependiendo del tipo de información que queramos almacenar.

En lenguajes como C, debemos declarar el tipo de variable que queremos usar. En lenguajes más recientes, como JavaScript o Python, no es necesario declarar el tipo de variable, ya que el lenguaje se encarga de inferirlo.

Existe entonces un compromiso: en algunos lenguajes, la persona que programa realiza mucho trabajo con mucho cuidado de planificar y diseñar el comportamiento de las variables, lo que permite que cuando el código corra, que lo haga de forma muy rápido. Por el otro lado, hay lenguajes que permiten que escribamos muy rápidamente código, sin preocuparnos por el tipo de variable, pero esto hace que cuando el código corra, lo haga de forma más lenta o ineficiente.

Cuando estamos haciendo tareas cotidianas en nuestro taller computacional, que correremos unas cuantas veces, o que no necesitamos rapidez, tendemos a prototipar en lenguajes modernos, como Python. Pero cuando estamos programando una máquina sonora, que necesita procesar audio en tiempo real, estar atenta a sensores, controlar actuadores, lo más probable es que necesitemos que esto ocurra muy rápido, por eso tenderemos a programar en lenguajes como C++.

Existe un paradigma importante: la cantidad de trabajo que se necesita para correr un programa debe repartirse entre dos entidades o momentos: 

\section{Booleanas}

Las variables booleanas son la partícula fundamental de las variables. Usan 1 bit de información, y pueden almacenar dos valores posibles:

\begin{table}
\begin{tabular}{c|c}
    \textbf{0} & \textbf{1} \\
    \hline
    False & Verdadero \\
    GND & $V_{cc}$ \\
\end{tabular}
\end{table}


\section{Enteras}

Las variables de tipo entera pueden almacenar números enteros, y usan más de un bit de información.

Existen varias variantes.

\section{Punto flotante}

Las variables de tipo punto flotante pueden almacenar números con parte decimal.

\section{Caracter}

Las variables de tipo caracter pueden almacenar un caracter, como una letra, un número.



\newpage

\end{document}

